\chapter{Capítulo II: Metodología}
\section{Metodología}
Para el desarrollo del sistema de gestión de nóminas para Café y
Vivero Verde Pradera, se empleará una combinación de metodologías ágiles,
específicamente el Desarrollo Incremental y Scrum. Esta combinación está
diseñada para adaptarse a la naturaleza dinámica del proyecto y facilitar
la colaboración continua con los usuarios finales.

\subsection{Fundamentación}

El Desarrollo Incremental permite la construcción del sistema en
múltiples ciclos o incrementos. Cada incremento agrega funcionalidad y
es plenamente operativo, lo que posibilita la entrega de una versión mejorada
del sistema con cada iteración. Esta metodología es particularmente efectiva
para obtener retroalimentación temprana y continua de los usuarios, permitiendo
ajustes que alinean el producto final con las necesidades específicas del cliente.

Por otro lado, Scrum se empleará para gestionar y planificar el proceso de
desarrollo. Scrum es un marco de trabajo que promueve la comunicación frecuente,
el trabajo en equipo, y la flexibilidad para adaptarse a cambios. Las ceremonias
de Scrum, como las reuniones diarias (Daily Scrum), las revisiones de sprint y
las planificaciones (Sprint Planning), facilitarán la coordinación y el
seguimiento constante del avance del proyecto.

\newpage

\subsection{Aplicación en el Proyecto}


En la aplicación práctica de estas metodologías en el proyecto:
\begin{itemize}
	\item \textbf{Planificación Scrum:} Se establecerán sprints de duración fija,
	      ya sea mensual o quincenal, durante los cuales se definirán objetivos
	      claros y alcanzables. Esto incluirá la definición de las funcionalidades
	      a desarrollar, pruebas y revisiones con los usuarios finales.
	\item \textbf{Desarrollo Incremental:} Cada sprint resultará en un
	      incremento funcional del sistema, que será evaluado por los usuarios finales.
	      La retroalimentación recogida será utilizada para mejorar o ajustar las
	      funcionalidades en los siguientes incrementos.
	\item \textbf{Cambios mínimos:} Se ajustará la duración de los sprints según
	      la complejidad de las funcionalidades y la capacidad del equipo de desarrollo,
	      buscando el equilibrio óptimo para una entrega eficiente y efectiva.
\end{itemize}





Con estas metodologías, el proyecto no solo apuntará a cumplir con los objetivos establecidos de manera eficaz, sino que también se adaptará a las necesidades cambiantes del entorno y los usuarios finales, asegurando la entrega de un sistema que verdaderamente responda a las expectativas y requerimientos de la empresa Café y Vivero Verde Pradera.

\subsection{Definición de estandares a utilizar}

En este apartado, describiremos los estándares seleccionados para cada una 
de las áreas clave del desarrollo del software, asegurando la consistencia, 
claridad y eficacia en la implementación y mantenimiento del proyecto.

\subsection{\textbf{Estándares de Documentación}}

Adoptaremos prácticas de documentación de código fuente, incluyendo comentarios detallados, 
encabezados de archivos y documentación exhaustiva de la API, para facilitar el entendimiento 
y la mantenibilidad del código. La documentación seguirá guías de estilo rigurosas que promueven
la claridad y un formato coherente en todos los documentos técnicos y manuales de usuario, 
asegurando que toda la información sea accesible y comprensible para todos los interesados.

Los comentarios deben realizarse de la siguiente manera:
\begin{itemize}
    \item \textbf{Bloque de comentarios:} Para explicar la funcionalidad de los métodos.
    \item \textbf{Línea de comentario:} Solo para la declaración de variables o en líneas de 
    código donde sea necesaria una especificación.
\end{itemize}

\subsection{\textbf{Estándares de Programación}}

Estableceremos convenciones de nombramiento para variables, clases y métodos en 
inglés que reflejen claramente su función. Utilizaremos un formato de indentación que 
facilite la lectura del código y mantendremos las líneas de código a una longitud razonable 
para evitar la complejidad visual. Emplearemos Git como sistema de control de versiones, con un 
esquema de nomenclatura de commits que incluye etiquetas como \texttt{[ADDED]}, \texttt{[BUG]}, 
\texttt{[FIXED]}, \texttt{[UPDATED]}, y \texttt{[FEATURE]}, acompañados de comentarios precisos 
que describan los cambios realizados. Además, incluiremos el uso de pull requests en GitHub para 
evitar conflictos en el desarrollo del software. Se prohibirá el uso de push forzados y la 
agregación de todos los archivos simultáneamente, a menos de que hayan sido modificados 
por la misma razón. Implementaremos revisiones de código entre pares y utilizaremos 
herramientas como VS Code Insiders con GitHub Copilot para realizar pruebas automáticas.

La estructura del proyecto se organizará de la siguiente manera: 
\begin{itemize}
    \item \textbf{Design:} Donde se incluirá el diseño de cada una de las pantallas a usar.
    \item \textbf{Docs:} Donde se almacenarán todos los documentos del sistema, como guías 
    de usuario y documentos de investigación.
    \item \textbf{Source:} Que contendrá el código fuente del sistema, distribuido en:
    \begin{itemize}
        \item \textbf{DB:} Que incluirá el diagrama de base de datos, el esquema de la 
        base de datos y el SQL de la base de datos.
        \item \textbf{API:} Donde se almacenará todo el código relacionado con la API.
        \item \textbf{ENV:} Que contendrá las variables de entorno.
        \item \textbf{Front-End:} Donde se incluirá todo lo relacionado con la parte frontal 
        del sistema.
        \item \textbf{Back-End:} Que contendrá todo lo relacionado con el backend del sistema.
    \end{itemize}
\end{itemize}

\subsection{\textbf{Estándares de Bases de Datos}}

Las definiciones de tablas y columnas seguirán un esquema de nomenclatura específico, 
donde las tablas comenzarán con \texttt{VPG\_} seguido del nombre de la tabla en mayúsculas, 
y las columnas se nombrarán en minúsculas siguiendo el formato \texttt{tablename\_atributo}. 
Las relaciones se nombrarán siguiendo el formato \texttt{VPG\_TABLENAME1\_TABLENAME2\_FK}. 
Adheriremos a las tres primeras reglas de normalización para garantizar una estructura de 
base de datos eficiente y optimizada tanto para consultas como para la integridad de los datos.

\subsection{\textbf{Estándares de Interfaz}}

Las interfaces seguirán una guía de estilo que incluye el uso de colores no muy 
llamativos basados en los colores de la empresa y la tipografía INTER para garantizar 
una presentación visual coherente y profesional. La disposición de los elementos estará a 
cargo del diseñador web, con un fuerte enfoque en la accesibilidad para cumplir con las 
directrices de WCAG. Además, fomentaremos la reutilización de componentes para maximizar la 
eficiencia del desarrollo. Las pruebas de interfaz se realizarán mediante evaluaciones 
entre compañeros y con feedback directo de los usuarios finales de la empresa, asegurando 
que las interfaces sean intuitivas y efectivas.

Estos estándares son fundamentales para el desarrollo exitoso del software, proporcionando un 
marco de trabajo claro y eficiente para todo el equipo del proyecto. Sin embargo, siempre habrá 
un espacio para revisar y adaptar las políticas conforme el proyecto crezca y cambie.